À medida que o número de objetos em órbita cresce, a estabilidade do ambiente operacional torna-se uma preocupação relevante.
Nas últimas décadas, acumulou-se uma quantidade expressiva de detritos espaciais associada a satélites, espaçonaves e estágios de foguetes, elevando o risco tanto para missões futuras quanto para operações em curso.
Esse conjunto inclui desde grandes componentes desativados até fragmentos oriundos de colisões e explosões em órbita.
Estima-se a existência de mais de 670.000 objetos com dimensões superiores a 1~cm em órbita terrestre \cite{esaSpaceDebris}, além de milhões de partículas menores que ainda podem representar perigo.

O cenário se torna pior por conta do aumento de constelações em órbita baixa, como a \emph{Starlink} e a \emph{OneWeb}, que buscam ampliar conectividade e serviços de comunicação.
Impactos em alta velocidade, conforme discutido em \cite{esaAboutSpaceDebris}, podem gerar grande quantidade de fragmentos, intensificando o problema e contribuindo para a síndrome de Kessler, na qual colisões em cascata aumentam progressivamente a densidade de detritos e podem comprometer determinadas faixas orbitais por períodos prolongados \cite{kessler1978}.

Além do risco de colisão, a concentração de detritos impõe consequências econômicas e operacionais.
Manobras de desvio elevam o consumo de propelente, aumentam a carga operacional e podem reduzir a vida útil de satélites.
Em contrapartida, a ausência de gestão adequada pode resultar em danos a ativos e interrupções de serviço, com efeitos diretos sobre a atividade orbital.

A gestão de detritos demanda coordenação internacional.
Medidas como captura e remoção de objetos por sistemas robóticos e a adoção de normas mais restritivas para desativação e descarte são caminhos para mitigar o problema.
Agências e organizações, como a \gls{esa} e a \gls{nasa}, também promovem diretrizes e iniciativas de mitigação \cite{esaDebrisMitigation2023}, incluindo o incentivo ao uso de materiais e projetos que favoreçam a desintegração completa durante a reentrada.
